\section{Research contributions}
\label{sec:contributions}

The main innovation of MouseBrainBench lies in its explicit separation between
model performance and the scientific claim derived from that performance. The
framework does not try to replace existing simulators or predictive benchmarks.
Instead, it provides a reproducible mechanism to decide whether a partial digital
model supports a predictive, reproducibility, structure--function, or
mechanistic-identifiability claim.

The contributions of the proposal are summarised in the following points:

\begin{itemize}
  \item The proposal introduces a claim-aware validation formulation that
  separates prediction, reproducibility, topology specificity, directed
  identifiability, local structure--function association, and computational
  cost. This fact allows the system to return not only a numerical result, but
  also an explicit set of admitted and rejected interpretations.

  \item A formal \acs{mis} is defined through non-interchangeable evidence
  blocks and fixed directional thresholds. The rule is evaluated with four
  known-truth systems that isolate common drive, topology without direction, and
  direction without topology specificity. These experiments validate the
  operational semantics of the claim gate, rather than providing an exhaustive
  Monte Carlo characterization of false-positive or false-negative rates.

  \item A cross-resource protocol is established by assigning Allen,
  Sensorium, Dynamic Sensorium, and \acs{microns} different roles in the
  validation process. Thus, heterogeneous datasets are not mixed as if they were
  repeated measurements of the same system. Allen acts as a real negative
  mechanistic control, Sensorium as a predictive control, and \acs{microns} as a
  local anatomical-functional case.

  \item A MICRONS validation case is provided using a discovery-fixed primary
  endpoint, distance- and degree-matched null models, false-discovery-rate
  correction, unit-cluster bootstrap stability intervals, and two non-overlapping hold-out
  windows. The aim is to recover a bounded and previously motivated local
  structure--function phenomenon, not to claim a new general wiring rule.

  \item The work includes a reproducible publication freeze that links numerical
  artifacts with allowed and blocked claims. In the current study, the framework
  admits an internally reproduced local observational structure--function association, while
  rejecting causal, behavioural, whole-brain, and state-of-the-art predictor
  claims. This produces a direct advantage over a conventional predictive
  evaluation, because strong reproducibility or prediction in Allen and Sensorium
  is not allowed to compensate for missing topology, direction, or causal
  evidence.
\end{itemize}

Notice that these contributions define the scope of the paper. MouseBrainBench
is a framework for evaluating the relationship between evidence and
interpretation. It does not establish that a single score is universally
sufficient for mechanistic neuroscience. The MIS thresholds and evidence blocks
must be specified according to the target claim before evaluation, and a failure
means that the analysed claim is not supported, not that the dataset or model is
without scientific interest.
